\chap{A boundary identification is data}\label{sec:boundary-lab}
The groupoid formulation can seem more elaborate than a set of equivalence classes until a change of presentation is examined closely. A two-view example is enough to show what is lost by freezing representatives instead of retaining their comparison maps. It also explains why the marked realization construction uses an identification at the boundary.

\subsection{Duplicate a view without adding information}
Let the original view category have one object and only its identity arrow. Replace it by a category with two objects, \(0\) and \(1\), and exactly one arrow between each ordered pair. These arrows are inverse isomorphisms. The two categories are equivalent: selecting object zero and forgetting the duplicate give inverse equivalences up to the evident isomorphism.

Use a constant group \(H\) as coefficients, with identity pullback homomorphisms. A cocycle on the duplicated category is specified by its label \(r\in H\) on the arrow \(0\to1\). Its reverse label is \(r^{-1}\), and identity labels are one. A gauge \((a_0,a_1)\) carries \(r\) to
\begin{equation}\label{eq:two-copy-gauge}
 r'=a_0ra_1^{-1}.
\end{equation}
Every \(r\) is isomorphic to the trivial cocycle, for example with gauge \((1,r)\). The automorphisms of \(r\) are the pairs
\[
 (a_0,r^{-1}a_0r),\qquad a_0\in H.
\]
Thus the unframed groupoid has one isomorphism class but retains a copy of \(H\) as the automorphism group of each object. This is compatible with the original one-view coefficient groupoid, which has one object and automorphism group \(H\).

\subsection{Freezing both representative gauges breaks that comparison}
Now require a based gauge to be the identity at every view. On the duplicated category the only permitted gauge is \((1,1)\). Equation~\eqref{eq:two-copy-gauge} then carries \(r\) only to itself. Distinct labels cannot be compared.

Restrict to the selected view zero. Every cocycle becomes the same one-view identity cocycle. The distinction between \(r\) and \(1\) has disappeared in the target, even though the fully framed source has no comparison arrow between them.

\begin{proposition}[Full representative framing is not invariant]\label{KS68}
For a nontrivial coefficient group, the restriction from the fully framed duplicated model to the fully framed single-view model need not be full, and hence need not be an equivalence. This failure occurs even though the underlying view categories are equivalent and the full prescribed domains are closed under isomorphism.
\end{proposition}
\begin{proof}
Choose \(r\ne1\). In the fully framed duplicated model there is no arrow from \(r\) to \(1\), because the only allowed gauge leaves \(r\) unchanged. After restriction to view zero, both become the same cocycle. There is an identity arrow between their images. That target arrow has no preimage, so the functor is not full. Fullness is necessary for an equivalence.
\end{proof}
\source{KS68}
The failure is not repaired by declaring the two objects equivalent as sets of labels. The issue is exactly the missing arrow. Nor does it contradict the theorem for marked relative realizations. That theorem does not freeze every arbitrary representative coordinate. It transports an actual boundary object together with an identification to it.

\subsection{The marked fiber keeps the comparison that matters}
In the same example, let \(R\) restrict an unframed duplicated cocycle to view zero. Fix the unique boundary object. A marked realization has two pieces of data:
\[
 (r,\phi),\qquad r\in H,\quad\phi\in H.
\]
The element \(\phi\) is the chosen identification of the restricted object with the fixed boundary. It is not discarded just because the boundary has only one object: that object still has many automorphisms.

A comparison from \((r,\phi)\) to \((r',\psi)\) is a gauge \((a_0,a_1)\) satisfying both the cocycle equation and the boundary equation
\[
 r'=a_0ra_1^{-1},\qquad \psi a_0=\phi.
\]
These equations determine it uniquely:
\begin{equation}\label{eq:marked-two-copy}
 a_0=\psi^{-1}\phi,\qquad a_1=(r')^{-1}a_0r.
\end{equation}
Substituting the formulas proves both equations without commuting any factors. Thus in this example there is exactly one identification-preserving arrow between any two marked realizations. The marked groupoid is nonempty and equivalent to the terminal groupoid, although its raw object set has many elements.

This is a worked finite model of the marked transport mechanism. The generic groupoid construction and its equivalence theorem are Proposition~\ref{KS38} and Theorem~\ref{KS39}. Formula~\eqref{eq:marked-two-copy} is checked in the independent finite suite for \(H=S_3\), retaining the order of multiplication. This finite check is separate from the general formal theorem.

\subsection{Four meanings of uniqueness}
The example separates four statements that are often compressed into the word unique.
\begin{center}
\begin{tabularx}{\linewidth}{@{}p{.28\linewidth}X@{}}
\toprule
Statement & What it retains\\
\midrule
One literal object & No different representative is present in the specified object set.\\
One isomorphism class & Different representatives may remain, as may their automorphisms.\\
A chosen realization & One object is selected, but no uniqueness follows from the choice.\\
One compatible comparison & The comparison respecting fixed identifications is unique, even when many representatives exist.\\
\bottomrule
\end{tabularx}
\end{center}
The unframed duplicated groupoid has one isomorphism class and nontrivial automorphisms. Its marked relative fiber has unique compatible comparisons. Its fully framed representative groupoid can instead have many mutually incomparable objects. These are three different constructions, not three descriptions of the same object with different terminology.

A descent theorem typically uses the last notion: two global realizations of the same specified local datum are related by a unique comparison respecting that datum. It does not force their underlying labels to be literally equal, or erase automorphisms after the local identifications have been forgotten.

\begin{insight}
Boundary data are not only values. In a symmetry-sensitive problem they include the identification with those values. Retaining that identification is what allows a relative realization problem to survive an equivalence of presentations.
\end{insight}
